\documentclass[11pt]{amsart}
\usepackage[T1]{fontenc}
\usepackage[utf8]{inputenc}
\usepackage{lmodern}
\usepackage[a4paper,margin=2.5cm]{geometry}
\usepackage{amsmath,amssymb,amsthm,mathtools}
\usepackage{booktabs,array,longtable}
\usepackage{enumitem}
\usepackage{xcolor}
\usepackage[colorlinks=true,linkcolor=blue!50!black,citecolor=green!40!black,urlcolor=blue!50!black]{hyperref}
\setlist{itemsep=2pt,topsep=3pt}
\setlength{\emergencystretch}{3em}

\theoremstyle{plain}
\newtheorem{theorem}{Theorem}[section]
\newtheorem{proposition}[theorem]{Proposition}
\newtheorem{lemma}[theorem]{Lemma}
\newtheorem{corollary}[theorem]{Corollary}
\newtheorem{conjecture}[theorem]{Conjecture}
\theoremstyle{definition}
\newtheorem{definition}[theorem]{Definition}
\newtheorem{hypothesis}[theorem]{Hypothesis}
\newtheorem{example}[theorem]{Example}
\theoremstyle{remark}
\newtheorem{remark}[theorem]{Remark}

% status tags
\newcommand{\status}[1]{\textup{\textsf{[#1]}}}
\newcommand{\stproved}{\status{proved}}
\newcommand{\stpmc}{\status{proved modulo cited theorems}}
\newcommand{\stcited}{\status{cited}}
\newcommand{\stsketch}{\status{sketched}}
\newcommand{\stexp}{\status{expected}}
\newcommand{\stconj}{\status{conjectural}}
\newcommand{\sthyp}{\status{hypothesis}}
\newcommand{\stopen}{\status{open}}
\newcommand{\stcomp}{\status{computed}}

\newcommand{\Q}{\mathbb{Q}}\newcommand{\Z}{\mathbb{Z}}\newcommand{\R}{\mathbb{R}}\newcommand{\C}{\mathbb{C}}
\newcommand{\F}{\mathbb{F}}\newcommand{\A}{\mathbb{A}}\newcommand{\cO}{\mathcal{O}}
\newcommand{\GL}{\mathrm{GL}}\newcommand{\SL}{\mathrm{SL}}\newcommand{\PGL}{\mathrm{PGL}}
\newcommand{\Nrd}{\mathrm{Nrd}}\newcommand{\Sym}{\mathrm{Sym}}
\newcommand{\JL}{\mathrm{JL}}\newcommand{\ord}{\mathrm{ord}}
\newcommand{\T}{\mathbb{T}}\newcommand{\m}{\mathfrak{m}}\newcommand{\cT}{\mathcal{T}}\newcommand{\cX}{\mathcal{X}}
\newcommand{\cM}{\mathcal{M}}\newcommand{\Frac}{\mathrm{Frac}}\newcommand{\tq}{\mathfrak{t}}\newcommand{\cc}{\mathfrak{c}}
\newcommand{\G}{\mathbb{G}}\newcommand{\Iw}{\mathrm{Iw}}\newcommand{\Hb}{\mathfrak{H}}
\newcommand{\fa}{\mathfrak{a}}\newcommand{\fb}{\mathfrak{b}}
\DeclareMathOperator{\Hom}{Hom}\DeclareMathOperator{\coker}{coker}\DeclareMathOperator{\End}{End}
\DeclareMathOperator{\Ann}{Ann}\DeclareMathOperator{\rk}{rk}\DeclareMathOperator{\Stab}{Stab}
\DeclareMathOperator{\tr}{tr}\DeclareMathOperator{\Jac}{Jac}

\title{Towards a simplicial approach to the Jacquet--Langlands correspondence over $\Q$}
\date{14 November 2026}

\begin{document}

\begin{abstract}
We describe a route towards a Jacquet--Langlands (JL) correspondence between a definite quaternion algebra $B$ over $\Q$ and $\GL_2$ over $\Q$ that is explicit, simplicial and integral.
Its intermediate objects are Bruhat--Tits trees at an auxiliary prime $q$ and the $S$-arithmetic groups acting on them.
All statements about the definite algebra and the ``indefinite'' stages are then finite combinatorics over class sets of $B$: quotient graphs of trees, a definite Ribet sequence obtained as a vertex collapse, a definite Ihara lemma, a weighted monodromy pairing and component group, and their ordinary $\Lambda$-adic limits.
These are proved from strong approximation and the connectivity of the tree.
The arithmetic geometry (Deligne--Rapoport models, Néron models, Grothendieck's monodromy, Čerednik--Drinfeld) that the classical arithmetic-geometric approach uses throughout is concentrated in a single, explicitly stated comparison hypothesis at the $\GL_2$ end.
As a first arithmetic output we record a $\Lambda$-adic integral JL \emph{defect formula} for $D=q$: the ratio of the congruence ideals of a Hida family and of its quaternionic transfer equals an old/new congruence ideal, which is computable from trees and amalgams, and also equals a Galois-theoretic Tamagawa ideal at $q$.
The document fixes notation, states what is proved, cited, sketched or conjectural, explains why the tree route is cheaper than the geometric route (for a Lean formalisation and for the readers of a final human-readable implementation), lists computations already carried out with the software HACS, and lists test cases for an implementation.
It is a seed for the detailed planning, implementation and formalisation of the route; it does not contain that plan.
\end{abstract}

\maketitle
\setcounter{tocdepth}{1}
\tableofcontents

%======================================================================
\section{Introduction}\label{sec:intro}
%======================================================================

\subsection{Aim}
The Jacquet--Langlands correspondence [JL] relates automorphic forms on the multiplicative group of a quaternion algebra to automorphic forms on $\GL_2$.
For a \emph{definite} quaternion algebra $B$ over $\Q$ of discriminant $D$, in its cohomological form it says: the Hecke eigensystems in degree-$0$ cohomology of the (finite) arithmetic quotients of $B^\times$, i.e.\ in Brandt modules, are the $D$-new Hecke eigensystems in degree-$1$ cohomology of modular curves.
Over $\C$ this is an equality of Hecke modules up to isomorphism; it is \emph{not} induced by a map of spaces, because the spaces involved have different dimensions (points versus Riemann surfaces).

The intention behind the programme seeded by this document is to find a \emph{homotopical} and \emph{explicitly simplicially realisable} approach to JL over $\Q$ in the definite case, with the following properties:
\begin{enumerate}[label=(\roman*)]
\item it is realised by Hecke-equivariant maps in \emph{both} directions (definite $\to\GL_2$ and $\GL_2\to$ definite), built from a zig-zag of spaces, simplicial orbifolds (in the sense of Moerdijk--Pronk [MP]) or chain complexes;
\item it is integral, at least over $\Z_{(p)}$ or $\Z_p$;
\item it specialises to the classical (semisimple) JL over $\C$;
\item mod $p^k$, in the limit over $p$-power levels, it specialises to Chenevier's $p$-adic JL [Ch];
\item it admits a (nearly) ordinary $\Lambda$-adic version that preserves congruences.
\end{enumerate}
Ichino and Prasanna [IP] show that integral and $p$-adic comparisons of this kind are delicate: they involve periods and $L$-values.
Any integral version must therefore produce, rather than ignore, a defect.
In the route described here the defect appears as a ratio of congruence ideals, and \S\ref{sec:defect} identifies it in the simplest case.

\subsection{The idea of the tree route}
Inverting an auxiliary prime $q$ replaces the arithmetic groups by $S$-arithmetic groups, which act on a product of a symmetric space and the Bruhat--Tits tree $\cT_q$ of $\PGL_2(\Q_q)$.
On the definite side the symmetric space is a point, so the norm-one $S$-arithmetic group of a maximal order of $B$ (for $q\nmid D$) acts on $\cT_q$ alone, with finite stabilisers.
Its quotient is a finite graph of finite groups whose vertices and edges are class sets of $B$ at two levels.
Thus a one-dimensional simplicial orbifold is built from zero-dimensional ones, and its first homology is the object that the arithmetic-geometric approach calls the character group of a Shimura curve (or modular curve) at a prime of bad reduction.

The route has three layers:
\begin{enumerate}[label=(\alph*)]
\item a \emph{combinatorial core} (\S\S\ref{sec:tree}--\ref{sec:lambda}): definite class sets, tree quotients, Ribet and Ihara statements, monodromy pairings, ordinary $\Lambda$-adic limits. Everything here is proved from strong approximation;
\item a \emph{comparison hypothesis} (\S\ref{sec:hypx}), $\mathrm{Hyp.X}$, which identifies the degree-zero Brandt module with the character group of the modular Jacobian at $q$, Hecke-equivariantly, with pairings and (optionally) Galois action. This is the only geometric input, and it is used only to \emph{interpret} definite objects as $\GL_2$ objects and to produce maps;
\item an \emph{arithmetic layer} (\S\ref{sec:defect}): congruence ideals, the old/new congruence module at $q$ and the integral JL defect, proved with Hida theory, a Gorenstein theorem of Hsieh and level lowering, independently of $\mathrm{Hyp.X}$.
\end{enumerate}

\subsection{What this document is}
It is a self-contained account of the mathematics of the route, as far as it has been developed, intended as the starting point for an independent exploration, implementation and formalisation.
We take the feasibility of the route for granted. We do not give a detailed work plan.
Every statement carries a status tag:
\stproved{} (a complete proof is given or is elementary and standard);
\stpmc{} (the proof is complete given the cited inputs, which are listed);
\stcited{} (taken from the literature);
\stsketch{} (an argument is indicated but not written out);
\stexp{} (a prediction, for instance of a computation);
\stconj{};
\sthyp{} (an explicit hypothesis of the route);
\stcomp{} (an exact computation that has been carried out).

\subsection{Structure}
\S\ref{sec:notation} fixes notation. \S\ref{sec:targets} recalls the classical, $p$-adic and $\Lambda$-adic targets.
\S\ref{sec:geom} recalls the arithmetic-geometric route and explains its cost.
\S\S\ref{sec:tree}--\ref{sec:lambda} develop the combinatorial core.
\S\ref{sec:hypx} states the comparison hypothesis and what it yields.
\S\ref{sec:defect} treats the integral defect.
\S\ref{sec:general} discusses several ramified primes.
\S\ref{sec:simplicial} describes the simplicial shape of the route.
\S\ref{sec:hacs} introduces HACS, lists the computations already done, and lists test cases.
\S\ref{sec:lean} discusses formalisation, and \S\ref{sec:open} collects open points.

%======================================================================
\section{Notation and conventions}\label{sec:notation}
%======================================================================

\subsection{Primes}
Throughout:
\begin{itemize}
\item $p\ge5$ is the \emph{coefficient} (or Hida) prime; $p$-power level towers are taken at $p$;
\item $q\ne p$ is the \emph{tree} prime: the auxiliary prime that is inverted, and at which the Bruhat--Tits tree lives;
\item $D$ is the discriminant of the definite quaternion algebra (a squarefree product of an odd number $\omega(D)$ of primes), with $p\nmid D$;
\item $M\ge1$ is a tame level with $(M,pqD)=1$ unless stated otherwise;
\item $\ell$ denotes an arbitrary prime (used for Hecke operators $T_\ell$), and $\omega(n)$ the number of prime factors of $n$.
\end{itemize}

\subsection{Adeles and quaternion algebras}
$\hat\Z=\prod_\ell\Z_\ell$, $\hat\Q=\hat\Z\otimes\Q$ (finite adeles). For a $\Q$-algebra $A$, $\hat A=A\otimes\hat\Q$ and $A_\ell=A\otimes\Q_\ell$.
$B=B_D$ is the definite quaternion algebra over $\Q$ ramified exactly at the primes dividing $D$ and at $\infty$; $\Nrd$ is the reduced norm, $B^1=\ker\Nrd$.
$\cO\subset B$ is a maximal order; for $\ell\nmid D$ we fix isomorphisms $\iota_\ell:\cO_\ell\xrightarrow{\sim}M_2(\Z_\ell)$.
For $\ell\mid D$, $B_\ell$ is a division algebra with maximal order $\cO_\ell$ and uniformiser $\Pi_\ell$.

\subsection{Levels and class sets}
For a compact open subgroup $U\subset\hat B^\times$,
\[X_B(U):=B^\times\backslash\hat B^\times/U\quad\text{(a finite set)},\qquad \Gamma_x:=B^\times\cap xUx^{-1}\quad(x\in\hat B^\times,\ \text{a finite group}).\]
Levels used (for $(N,D)=1$; matrices read through $\iota$):
\begin{itemize}
\item $U_0(N)$, $U_1(N)\subset\hat\cO^\times$: elements congruent to $\left(\begin{smallmatrix}*&*\\0&*\end{smallmatrix}\right)$, resp.\ $\left(\begin{smallmatrix}*&*\\0&1\end{smallmatrix}\right)$, modulo $N$;
\item the $p$-tower $U_s:=U_1(Mp^s)$ ($s\ge0$), $S_s:=X_B(U_s)$, and $\Delta_s:=(1+p\Z_p)/(1+p^s\Z_p)$ acting by diamond operators;
\item for $q\nmid D$ and $U$ with $U_q=\GL_2(\Z_q)$: $U_0(q)U:=U\cap U_0(q)$, i.e.\ $U_q$ replaced by the Iwahori subgroup $\Iw_q=\{\left(\begin{smallmatrix}a&b\\c&d\end{smallmatrix}\right)\in\GL_2(\Z_q):q\mid c\}$.
\end{itemize}
All these satisfy $\Nrd(U)=\hat\Z^\times$ \stproved{} (determinant is surjective on each local factor).

\subsection{Coefficients and Brandt modules}
$V_k:=\Sym^{k-2}(\Z_p^2)$ ($k\ge2$ even) with $\GL_2(\Z_p)$ acting on homogeneous polynomials of degree $k-2$; we write $n=k-2$ when the symmetric power degree is meant.
For a finitely generated $\Z_p$-module $V$ with continuous action of $U_p\subseteq\GL_2(\Z_p)$, the \emph{Brandt module} is
\[M(U;V):=\{f:\hat B^\times\to V:\ f(\gamma xu)=u_p^{-1}f(x)\ (\gamma\in B^\times,\ u\in U)\}\ \cong\ \bigoplus_{[x]\in X_B(U)}V^{\,x_p^{-1}\Gamma_xx_p}.\]
For $V=\Z$ (or $\Z_p$) we write $\Z[X]$ for the free module on a finite set $X$ (homological, pushforward) and $\Z^X$ for functions (cohomological, pullback); $\Z[X]^0$ is the degree-zero submodule.
When $p\nmid|\Gamma_x|$ the two are identified over $\Z_p$ by the norm map.
The weights $w_x:=|\Gamma_x/(\Gamma_x\cap\{\pm1\})|$ define the Brandt (mass) pairing $\langle[x],[y]\rangle=w_x\delta_{xy}$; whether one uses $|\Gamma_x|$ or $|\Gamma_x/\pm1|$ changes it by a factor $2$, and every computation must say which.

\subsection{Hecke operators}
For $\ell\nmid Dpq$ and $g\in\hat B^\times$ supported at $\ell$, write $[UgU]=\coprod_ig_iU$ and $(T_gf)(x)=\sum_if(xg_i)$.
This gives $T_\ell$ ($g=\mathrm{diag}(\ell,1)$), $S_\ell$ ($g=\mathrm{diag}(\ell,\ell)$), $U_\ell$ for $\ell\mid M$, diamond operators $\langle d\rangle$ (right translation by $\mathrm{diag}(1,d)$ at the primes dividing $Mp$), and $U_p=[U\,\mathrm{diag}(1,p)\,U]$ with the semigroup action on $V$.
Of the two natural normalisations of $U_p$ only one commutes with pushforward along $S_{s+1}\to S_s$; computations must fix it.
For $\ell\mid D$, $\Pi_\ell$ acts by right translation (the quaternionic $U_\ell$).
All operators supported away from $q$ commute with right translation by $B_q^\times$.

On the $\GL_2$ side, $\Gamma_0(N),\Gamma_1(N)\subset\SL_2(\Z)$ are the usual congruence subgroups, $X_0(N)$, $X_1(N)$ the modular curves over $\Q$, $J_0(N)$, $J_1(N)$ their Jacobians, $S_k(\Gamma)$ cusp forms, and $H^1_{\rm par}(\Gamma,V)$ parabolic cohomology. Hecke operators are normalised classically: $T_\ell$, $U_\ell$ act on a normalised eigenform by its $\ell$-th Fourier coefficient.

\subsection{Hida theory}
$\Lambda:=\cO_E[[1+p\Z_p]]$ for the ring of integers $\cO_E$ of a finite extension $E/\Q_p$ (large enough), with uniformiser $\varpi$; $P_k\subset\Lambda$ is the kernel of $[u]\mapsto u^{k-2}$ (weight $k$, trivial character).
$e=\lim_nU_p^{n!}$ is Hida's ordinary projector; it exists and is exact on finite $\Z/p^m$-modules and on profinite $\Z_p[[U_p]]$-modules with continuous $U_p$-action.
For a $\Lambda$-algebra $S$ and an algebra map $\mu:S\to A$ to a domain $A$, the \emph{congruence ideal} is $\eta(S,\mu):=\mu(\Ann_S\ker\mu)\subseteq A$.

\subsection{Weighted graphs}
A weighted graph $\G$ is a finite graph (multiple edges allowed) with edge weights $w_e\in\Z_{\ge1}$. Its chain complex is $\Z[E]\xrightarrow{\partial}\Z[V]$; put $X(\G):=H_1(\G)=\ker\partial$,
\[\langle e,e'\rangle_w:=w_e\delta_{ee'},\qquad \Phi^{\rm comb}(\G):=\coker\big(X(\G)\to\Hom(X(\G),\Z),\ x\mapsto\langle x,\cdot\rangle_w\big).\]
We call $\langle,\rangle_w$ the \emph{combinatorial monodromy pairing} and $\Phi^{\rm comb}$ the \emph{combinatorial component group}.

%======================================================================
\section{The targets}\label{sec:targets}
%======================================================================

\subsection{Classical JL over $\C$}
\begin{theorem}[{[JL], Eichler, Shimizu; cohomological form}]\label{thm:JLclassical}
Let $k\ge2$ be even and $(M,D)=1$. There is an isomorphism of modules over the Hecke algebra generated by the $T_\ell$, $\ell\nmid DM$,
\[M(U_0(M);V_k\otimes\C)\ \cong\ S_k(\Gamma_0(DM))^{D\text{-new}}\ \oplus\ \begin{cases}\C\ (\text{Eisenstein, the constants})&k=2,\\0&k>2,\end{cases}\]
and by Eichler--Shimura $H^1_{\rm par}(\Gamma_0(DM),V_k\otimes\C)\cong S_k\oplus\overline{S_k}$. \stcited
\end{theorem}
So over $\C$ a Hecke-equivariant map $H^0(\text{definite})\to H^1(\text{modular curve})$ exists, but it is only defined up to the choice of an isomorphism on each eigenspace and of a sign $\pm$; nothing canonical or integral is asserted.

\subsection{Integral and $p$-adic targets}
\emph{Integrally}, the Hecke algebras agree (for instance, the $\Lambda$-adic ordinary Hecke algebra of $\GL_2$, quotiented to its $D$-new part, is isomorphic to the definite one [LV], [Hs]), but the lattices do not: the congruence ideals on the two sides differ (\S\ref{sec:defect}), and their ratio is governed by local Tamagawa factors and periods [PW], [KO], [IP].
\emph{$p$-adically}, Chenevier [Ch] constructs a closed immersion of the eigencurve of the definite algebra into the eigencurve of $\GL_2$ (of matching tame level), with image the $D$-new part, interpolating classical JL \stcited. Its proof works with $p$-adic families of automorphic forms (overconvergent forms, Fredholm theory); it does not produce maps of integral cohomology or of spaces.
The aim (iv) of \S\ref{sec:intro} is a version that comes from the integral simplicial objects of this document by reduction mod $p^k$ and passage to the limit.

\subsection{$\Lambda$-adic target}
On the ordinary locus Hida theory replaces the eigencurve by finite free $\Lambda$-algebras and modules. The route of this document produces, on the definite side, $\Lambda$-adic modules and exact sequences (\S\ref{sec:lambda}), and, under the comparison hypothesis, $\Lambda$-adic maps to and from Hida's ordinary Tate module (\S\ref{sec:hypx}).
Over $\Q$ the nearly ordinary theory is the ordinary theory tensored with cyclotomic twists, so nothing beyond the ordinary case is needed \stsketch.

%======================================================================
\section{Context: the arithmetic-geometric route and its cost}\label{sec:geom}
%======================================================================

\subsection{The geometric route}\label{ss:geomroute}
The classical integral link between Brandt modules and modular forms is Deligne--Rapoport (DR) and Ribet [DR], [Ri90], [RS].
Let $D=q$, $N=qM'$ with $(q,M')=1$, and $J=J(\Gamma_1(M')\cap\Gamma_0(q))$.
\begin{enumerate}[label=(G\arabic*)]
\item $X(\Gamma_1(M')\cap\Gamma_0(q))$ has semistable reduction at $q$. Its special fibre consists of two copies of $X_1(M')_{\F_q}$ crossing transversally at the supersingular points, and the supersingular points are in bijection with a class set of $B_q$ (Deuring, Serre). \stcited{} [DR], [RS, \S3]
\item The toric part of the Néron model of $J$ at $q$ has character group $X_q(J)\cong\Z[X_{B_q}(U)]^0$, Hecke-equivariantly; Frobenius acts on the torus as $qU_q$. \stcited{} [RS, (3.1), Lemma~3.17], [PW, Prop.~6.5] (at $\Gamma_0$-level)
\item Grothendieck's monodromy pairing on $X_q(J)$ is the weighted pairing $w_e\delta_{ee'}$ (node thicknesses), and its cokernel is the component group $\Phi_q(J)$. \stcited{} [SGA7, IX], [Ill]
\item The inertia filtration $T_t\subset T_pJ^{I_q}\subset T_pJ$ has graded pieces $\Hom(X_q(J),\Z_p(1))$, $T_pJ_1(M')^{\oplus2}$, $X_q(J)\otimes\Z_p$. Hence there are Hecke-equivariant maps $\Hom(X_q,\Z_p(1))\hookrightarrow T_pJ$ and $T_pJ\twoheadrightarrow X_q\otimes\Z_p$: an \emph{integral cohomological JL in both directions}, depending on a choice of embedding $\iota:\bar\Q\hookrightarrow\bar\Q_q$. \stcited{} [SGA7, IX]
\item For $\omega(D)\ge3$ one passes through Shimura curves: Čerednik--Drinfeld $q$-adic uniformisation [Če], [Dr], [BC] describes their reduction at $q\mid D$ by tree quotients for a definite algebra, and Ribet's exact sequence of character groups [Ri90] links character groups at two primes. \stcited
\end{enumerate}
At $p$-power level and in Hida's limit, (G1)--(G4) give a $\Lambda$-adic monodromy filtration of Hida's ordinary Tate module with $\Lambda$-adic definite pieces \stsketch.

\subsection{Why the geometric route is expensive}\label{ss:cost}
\emph{Prerequisites for a reader.}
(G1)--(G5) use moduli of elliptic curves over $\Z_{(q)}$ and their regular models, supersingular reduction and Deuring's correspondence, Néron models and their component groups, semistable reduction of Jacobians, Grothendieck's monodromy theorem with its pairing (SGA~7), $\ell$-adic or étale cohomology (or at least Tate modules with Galois action), and, for several ramified primes, Shimura curves and $q$-adic uniformisation by Drinfeld's upper half-plane.
A human-readable final account of an implementation along the geometric route would therefore assume a full graduate course in arithmetic geometry.
The tree route needs: quaternion algebras over $\Q$ (orders, class sets, strong approximation), Bruhat--Tits trees, cohomology of groups acting on trees, elementary commutative algebra, and Hida theory; Galois representations enter only in \S\ref{sec:defect}, as cited black boxes.

\emph{Formalisation.}
In Lean~4 with Mathlib, none of the geometric objects above (moduli of elliptic curves over $\Z_{(q)}$, Néron models, Grothendieck's monodromy pairing, Shimura curves, $q$-adic uniformisation) is available in a form that can be reused for this purpose.
Large existing formalisations around Fermat's Last Theorem contain fragments: a Deligne--Rapoport-type model of $X_0(N_0q)$ with its special fibre, the Deuring bijection with Hecke compatibility at $\Gamma_0$-type level, the toric part of $J_0$ as $\Hom(X,\bar\F_q^\times)$, as well as the tree, the amalgam $\SL_2(\Z[1/q])$ and strong approximation in the shape needed here.
The geometric results sit in import closures of the order of $10^5$--$10^6$ lines, on other Mathlib versions, and they do not cover node thicknesses, $\Gamma_1(Mp^s)$-level, the $p$-tower or Čerednik--Drinfeld.
The pure, combinatorial pieces (tree, amalgam, strong approximation, coset graphs) are small.
A formalisation along the geometric route would therefore either import a very large, moving dependency or carry many deep geometric statements as hypotheses at every level.
The tree route carries \emph{one} such hypothesis, needed only for interpretation and for maps (\S\ref{sec:hypx}).

\emph{Canonicity.}
The geometric maps of (G4) depend on $\iota$ and are not induced by maps of Betti cohomology. The tree route does not recover this canonicity, but it loses nothing that the geometric route had in Betti terms.

%======================================================================
\section{Trees, $S$-arithmetic quaternion groups and quotient graphs}\label{sec:tree}
%======================================================================

\subsection{The Bruhat--Tits tree}
$\cT_q$ has as vertices the homothety classes $[L]$ of $\Z_q$-lattices $L\subset\Q_q^2$; $[L]$ and $[L']$ are adjacent iff they have representatives with $qL\subsetneq L'\subsetneq L$. It is a $(q+1)$-regular tree [Se, II.1.1] \stcited; a Lean formalisation exists [LM].
The type $\tau[g\Z_q^2]:=\ord_q\det g\bmod2$ is well defined; $g\in\GL_2(\Q_q)$ preserves types iff $\ord_q\det g$ is even. Put $v_0=[\Z_q^2]$, $\alpha=\mathrm{diag}(1,q)$, $v_1=\alpha v_0$, $e_0=\{v_0,v_1\}$.

\begin{lemma}\label{lem:tree}
$\SL_2(\Q_q)$ acts transitively on type-$0$ vertices, on type-$1$ vertices and on edges, with stabilisers $\SL_2(\Z_q)$, $\SL_2(\Q_q)\cap\alpha\GL_2(\Z_q)\alpha^{-1}$ and $\SL_2(\Q_q)\cap\Iw_q$. Moreover $\GL_2(\Z_q)\cap\alpha\GL_2(\Z_q)\alpha^{-1}=\Iw_q$. \stproved
\end{lemma}
\begin{proof}
If $\ord\det g=2m$ and $\det g=q^{2m}u$, then $g'=q^{-m}g\,\mathrm{diag}(u^{-1},1)\in\SL_2(\Q_q)$ and $g'v_0=gv_0$; likewise for type $1$ with $v_1$. Every edge has one vertex of each type, and $\SL_2(\Z_q)$ acts on the $q+1$ neighbours of $v_0$ through $\SL_2(\F_q)$ on $\mathbb P^1(\F_q)$, transitively. An element of $\SL_2$ stabilising $e_0$ fixes both vertices. Finally $\alpha^{-1}\left(\begin{smallmatrix}a&b\\c&d\end{smallmatrix}\right)\alpha=\left(\begin{smallmatrix}a&qb\\c/q&d\end{smallmatrix}\right)$ is integral iff $q\mid c$.
\end{proof}

\subsection{The $S$-arithmetic group}
Assume $q\nmid D$, $U_q=\iota_q^{-1}\GL_2(\Z_q)$ and $\Nrd(U)=\hat\Z^\times$. Put
\[\Gamma_U:=\{\gamma\in B^\times:\ \Nrd\gamma=1,\ \gamma\in U_\ell\text{ for all }\ell\ne q\}=\cO_U[\tfrac1q]^1,\]
acting on $\cT_q$ through $\iota_q$, type-preservingly.

\begin{lemma}[stabilisers]\label{lem:stab}
(a) $\Nrd(\Gamma_x)=1$ for every $x$. (b) If $U\subseteq U_1(p^s)$ with $s\ge1$, $p\ge5$, then all $\Gamma_x$ are trivial and $\Delta_s$ acts freely on $X_B(U)$. (c) For the levels of \S\ref{sec:notation} and $p\ge5$, $p\nmid|\Gamma_x|$. \stproved
\end{lemma}
\begin{proof}
(a) $\Nrd\gamma\in\Q_{>0}\cap\hat\Z^\times=\{1\}$. (b) A non-trivial $\gamma\in\Gamma_x$ has finite order and characteristic polynomial $X^2-tX+1$, $t\in\{-2,-1,0,1\}$. At $p$ it is conjugate into $U_1(p)$, so $1$ is an eigenvalue mod $p$ and $p\mid2-t\in\{4,3,2,1\}$, impossible. If $\gamma xk=x\,\mathrm{diag}(1,d)$ then $\gamma=1$ and $d\in1+p^s\Z_p$. (c) Element orders in a quaternion order over $\Z$ lie in $\{1,2,3,4,6\}$.
\end{proof}

\begin{proposition}[quotient graph]\label{prop:quot}
\begin{enumerate}[label=(\alph*)]
\item There are bijections $\Gamma_U\backslash\{\text{type-}j\text{ vertices}\}\cong X_B(U)$ ($j=0,1$) and $\Gamma_U\backslash\{\text{edges}\}\cong X_B(U_0(q)U)$, under which the endpoint maps are the two degeneracies
\[s_0[x]=[x],\qquad s_1[x]=[x\alpha]\qquad(\alpha=\iota_q^{-1}\mathrm{diag}(1,q)).\]
\item The stabiliser in $\Gamma_U$ of the cell corresponding to $[x]$ is $\Gamma_x$ (at the relevant level).
\item The bijections are equivariant for all Hecke operators supported away from $q$, and $s_0,s_1$ commute with them.
\end{enumerate}
\stpmc: strong approximation for $B^1$ with respect to $S=\{\infty,q\}$ ($B_q$ is split, so $B^1$ is dense in $\hat B^{1,(q)}$) [Vo, Ch.~28].
\end{proposition}
\begin{proof}
Let $U'\in\{U,U_0(q)U,\alpha U\alpha^{-1}\}$; all have $U'^{(q)}=U^{(q)}$ and $\Nrd U'=\hat\Z^\times$.
Since $\Nrd(B^\times)=\Q_{>0}$ and $\Q_{>0}\hat\Z^\times=\hat\Q^\times$, every class of $X_B(U')$ has a representative of norm in $\hat\Z^\times$, and two such differ by $\gamma$ of norm $1$; hence $X_B(U')=B^1\backslash\hat B^1/U'^1$.
Strong approximation gives $\hat B^1=B^1(B^1_q\times U'^{1,(q)})$, so $X_B(U')=\Gamma_U\backslash B^1_q/U'^1_q$ (injectivity: if $\beta gu=g'$ then $\beta\in\Gamma_U$).
By Lemma~\ref{lem:tree}, $B^1_q/U'^1_q$ is the set of type-$0$ vertices, type-$1$ vertices or edges. The edge $ge_0$ has endpoints $gv_0$, $g\alpha v_0$, which gives $s_0,s_1$; $s_1$ is well defined because $\alpha^{-1}\Iw_q\alpha\subseteq\GL_2(\Z_q)$.
(b) is the computation of stabilisers in $B^1_q$; (c) holds because Hecke operators away from $q$ act by right translation at places $\ne q$.
\end{proof}

\begin{remark}[why norm one]\label{rem:normone}
With the full group $\cO_U[\frac1q]^\times$ the central element $q$ acts trivially on $\cT_q$, and at $U_1(p^s)$-level the central idele $q^{(q)}$ acts on $X_B(U)$ as a diamond operator, so that vertex orbits become $X_B(U)/\langle q\rangle$. The norm-one group gives two copies of $X_B(U)$, which is what the Ribet sequence needs.
\end{remark}

\subsection{Cohomology with coefficients}
\begin{proposition}\label{prop:coh}
Let $p\ge5$ and let $V$ be a finitely generated $\Z_p$-module with continuous $U_p$-action, $\Gamma_U$ acting through $\Gamma_U\subset U_p$. Then $H^*(\Gamma_U,V)$ is the cohomology of the two-term complex
\[0\to M(U;V)^{\oplus2}\xrightarrow{(f_0,f_1)\mapsto s_1^*f_1-s_0^*f_0}M(U_0(q)U;V)\to0,\]
and $H_*(\Gamma_U,V)$ the homology of the transpose $\partial=(-s_{0*},s_{1*})$. In particular $H^i=H_i=0$ for $i\ge2$, and $H_0(\Gamma_U,V)=V_{U^1_p}$ (coinvariants under the closure $U_p^1$ of $\Gamma_U$). For $V=V_k$, $k>2$, $V^{\Gamma_U}=0$ and $H_0$ is finite. \stpmc: the equivariant spectral sequence for a group acting on a tree without inversions [Br, VII.7], [Se].
\end{proposition}
\begin{proof}
$\cT_q$ is contractible, the action has no inversions and the stabilisers have order prime to $p$ (Lemma~\ref{lem:stab}), so the spectral sequence degenerates to invariants of stabilisers, which Prop.~\ref{prop:quot} identifies with Brandt modules. $H_0$: coinvariants are stable under the closure, which is $U_p^1$ by strong approximation. For $k>2$, $V_k\otimes\Q_p$ is an irreducible non-trivial representation of a Zariski-dense subgroup of $\SL_2$.
\end{proof}

\begin{proposition}[harmonic cochains]\label{prop:harm}
The $\Gamma_U$-invariant $V$-valued harmonic cochains on $\cT_q$ (functions $c$ on oriented edges with $c(\bar e)=-c(e)$ and $\sum_{s(e)=v}c(e)=0$) are $\ker\big(M(U_0(q)U;V)\xrightarrow{(s_{0*},s_{1*})}M(U;V)^2\big)\cong H_1(\Gamma_U,V)$, the $q$-new part of the Brandt module of level $U_0(q)U$. \stproved
\end{proposition}
Here $V_k$ is an \emph{integral} coefficient system for the $S$-arithmetic group because the tree lives at $q\ne p$: elements of $\cO[\frac1q]^\times$ are $p$-integral. (This differs from the Teitelbaum--Darmon setting, where the tree is at $p$.)

\subsection{Indefinite stages as tree quotients}
\begin{definition}\label{def:indef}
For $q\nmid D$ put $\cX^{(q)}_{Dq}(U;V):=H_1(\Gamma_U,V)$, the \emph{tree model at $q$} of the indefinite discriminant $Dq$.
\end{definition}
By Čerednik--Drinfeld, for $V=\Z$ and $\Gamma_0$-type $U$ this is the character group at $q$ of the Jacobian of the Shimura curve of discriminant $Dq$ \stcited{} [Če], [Dr], [BC]. The tree route never uses this identification. An indefinite discriminant $d$ has $\omega(d)$ tree models (one for each $q\mid d$, with definite algebra $B_{d/q}$), and nothing combinatorial relates them (\S\ref{sec:general}).

%======================================================================
\section{Graph models, the Ribet sequence and the Ihara lemma}\label{sec:graph}
%======================================================================

\subsection{Weighted graphs}
For a semistable curve with dual graph $\G$ and node thicknesses $w_e$, Grothendieck identifies $X(\G)$ with the character group of the toric part of the Néron model, $\langle,\rangle_w$ with the monodromy pairing and $\Phi^{\rm comb}$ with the component group [SGA7, IX], [Ill] \stcited. In the tree route this dictionary is \emph{not} used: $X$, $\langle,\rangle_w$ and $\Phi^{\rm comb}$ are definitions.

\begin{proposition}[two-vertex graphs]\label{prop:phi}
If $\G$ has two vertices and all edges $E$ join them, then $X(\G)=\Z[E]^0$ and
\[|\Phi^{\rm comb}(\G)|=\Big(\prod_ew_e\Big)\sum_e\frac1{w_e}.\]
If all $w_e=1$, $\Phi^{\rm comb}\cong\Z/|E|$; in general it need not be cyclic ($w=(2,2,2)$ gives $\Z/2\times\Z/6$). \stproved
\end{proposition}
\begin{proof}
$\Hom(\Z[E]^0,\Z)=\Z^E/\Z\mathbf 1$ and the map is $x\mapsto(w_ex_e)\bmod\mathbf1$, so $\Phi^{\rm comb}=\Z^E/(W\Z[E]^0+\Z\mathbf1)$. Its order is the weighted spanning-tree count $\sum_e\prod_{e'\ne e}w_{e'}$.
\end{proof}

\begin{definition}[DR graph]\label{def:DR}
For $B=B_q$ and $U=U_1(Mp^s)$ or $U_0(M)$, $\G^{\rm DR}_q(U)$ is the two-vertex weighted graph with edge set $X_{B_q}(U)$ and $w_e=|\Gamma_e/(\Gamma_e\cap\{\pm1\})|$. Put $X^{\rm comb}_q(U):=\Z[X_{B_q}(U)]^0$, and for $k>2$, $X^{\rm comb}_q(U;V_k):=M(U;V_k)$.
\end{definition}
For $s\ge1$ all $w_e=1$ (Lemma~\ref{lem:stab}) and $\Phi^{\rm comb}\cong\Z/|S_s|$.
At level $1$ and $q\ge5$, Eichler's mass formula $\sum_e|\cO_e^\times|^{-1}=\frac{q-1}{24}$ and the classes with extra units (one with $w=2$ iff $q\equiv3\bmod4$, one with $w=3$ iff $q\equiv2\bmod3$) give $|\Phi^{\rm comb}|=\mathrm{num}\frac{q-1}{12}$, the order of Mazur's component group of $J_0(q)$ [Mz] \stpmc.

\begin{example}\label{ex:mazur}
$q=11$: $w=(3,2)$, $\Phi^{\rm comb}=\Z/5$. $q=23$: $w=(1,2,3)$, $|\Phi^{\rm comb}|=11$. $q=37$: $w=(1,1,1)$, $\Phi^{\rm comb}=\Z/3$. All equal $\mathrm{num}\frac{q-1}{12}$. This is a purely definite check of the pairing part of the comparison hypothesis at an \emph{Eisenstein} maximal ideal.
\end{example}

\subsection{The definite Ribet sequence}
\begin{theorem}[definite Ribet sequence as a vertex collapse]\label{thm:ribet}
Let $q\nmid D$, $U$ as in Prop.~\ref{prop:quot}, and $\G:=\Gamma_U\backslash\backslash\cT_q$ the quotient graph of groups, $|\G|$ its underlying graph.
\begin{enumerate}[label=(\alph*)]
\item There is an exact sequence over $\Z$, equivariant for all Hecke operators away from $q$:
\[0\to H_1(|\G|,\Z)\to\Z[X_B(U_0(q)U)]^0\xrightarrow{\ \alpha=(-s_{0*},\,s_{1*})\ }\big(\Z[X_B(U)]^0\big)^{2}\to0.\]
It is the long exact sequence of the collapse $|\G|\to\G^{\rm DR}$ onto the graph with the same edges and two vertices.
\item For $p\ge5$ and any $V$ as in Prop.~\ref{prop:coh}:
\[0\to H_1(\Gamma_U,V)\to M(U_0(q)U;V)\xrightarrow{\alpha}M(U;V)^{2}\to H_0(\Gamma_U,V)\to0.\]
\end{enumerate}
\stproved{} (given Prop.~\ref{prop:quot}).
\end{theorem}
\begin{proof}
(a) The collapse $c$ is the identity on edges and sends type-$j$ vertices to $j$. Its kernel complex is $\Z[V_0]^0\oplus\Z[V_1]^0$ in degree $0$. Since $|\G|$ is connected (a quotient of the tree) and $\G^{\rm DR}$ is connected, the long exact sequence reads $0\to H_1(|\G|)\to\Z[E]^0\to\Z[V_0]^0\oplus\Z[V_1]^0\to0$; identify $E$, $V_j$ by Prop.~\ref{prop:quot}. (b) is the homology of the complex of Prop.~\ref{prop:coh}.
\end{proof}

\begin{corollary}[definite Ihara lemma]\label{cor:ihara}
(i) $\alpha$ in (a) is surjective over $\Z$. (ii) For every prime $\ell$, the kernel of $\alpha^*=(-s_0^*,s_1^*)\bmod\ell$ on $(\F_\ell^{X_B(U)})^2$ is the line of constants $(\mathbf1,\mathbf1)$. (iii) With coefficients, $\coker\alpha=H_0(\Gamma_U,V)=V_{U^1_p}$; a Hecke operator $T$ supported at $\ell\nmid Dpq$ acts on it as $\deg(T)$ times the class of an element of $U_p$ of determinant $\ell^{\pm1}$. Hence $H_0(\Gamma_U,V)$ is Eisenstein and vanishes after localisation at a non-Eisenstein maximal ideal.
(i), (ii) \stproved{} (connectivity of $\cT_q$); (iii) identification \stproved, Hecke action \stsketch.
\end{corollary}
The definite Ihara lemma via strong approximation is classical [DT]; what is new here is only the integral form at $U_1(Mp^s)$-level with the complete proof above, and its use as the definition of the indefinite stage.

\begin{remark}
Ribet's sequence for $\GL_2$ character groups [Ri90] is Thm.~\ref{thm:ribet} plus two identifications that the tree route does not prove: $H_1(|\G|)=X_q(J^{Dq})$ (Čerednik--Drinfeld) and $\Z[X_{B}(U_0(q)U)]^0=X_r(J)$ at another prime $r$ (the comparison hypothesis of \S\ref{sec:hypx}).
\end{remark}

\subsection{The $\GL_2$ side: the amalgam}\label{ss:amalgam}
$\Gamma_0(M)[\frac1q]=\Gamma_0(M)*_{\Gamma_0(Mq)}\alpha'^{-1}\Gamma_0(M)\alpha'$ with $\alpha'=\mathrm{diag}(q,1)$ [Se, II.1.4] \stcited. It acts on $\Hb\times\cT_q$; its quotient is a graph of \emph{spaces} $[Y_0(M)\sqcup Y_0(M)\leftleftarrows Y_0(Mq)]$. Mayer--Vietoris gives
\[0\to H^1(\Gamma_0(M)[\tfrac1q],V)\to H^1(\Gamma_0(M),V)^2\xrightarrow{\Psi}H^1(\Gamma_0(Mq),V)\to H^2(\Gamma_0(M)[\tfrac1q],V)\to0\]
for $V=\Sym^{k-2}$ over $\Z_{(l)}$, $l>k$, $l\ne q$. The $q$-new part moves up one degree, into $H^2$, so the amalgam alone does not see $B_q$. Its use in the route is as the \emph{Betti computation of the $\GL_2$ old/new congruence module at $q$}: the torsion of $\coker\Psi$ (an Ihara lemma for $\GL_2$, Eisenstein only \stexp) and the Ribet level-raising module $\coker(\Psi^*\Psi)$ [Ri84], whose localisation at a cuspidal maximal ideal computes the fixed-weight old/new congruence module $\cc_q$ of \S\ref{sec:defect} \stsketch. Both have been computed (\S\ref{sec:hacs}).

%======================================================================
\section{Ordinary and $\Lambda$-adic passage}\label{sec:lambda}
%======================================================================

\begin{lemma}[$e$ kills degree, vertex and component terms]\label{lem:ekills}
Let $s\ge1$, $V=\Z_p$. (a) $e$ kills $\Z_p$ with $U_p$ acting through its degree $p$; hence $e\Z_p[S_s]^0=e\Z_p[S_s]$ and $e$ kills $H_0(\Gamma_U,\Z_p)$. (b) $e$ kills the vertex term of $\G^{\rm DR}$, so $eX^{\rm comb}_q=e\Z_p[S_s]$. (c) $e\big(\Phi^{\rm comb}\otimes\Z_p\big)=0$, and the pairing restricts to a perfect pairing $e\Z_p[S_s]\times e^*\Z_p[S_s]\to\Z_p$. \stproved
\end{lemma}
\begin{proof}
The degree intertwines $U_p$ with multiplication by $p$, which is topologically nilpotent; $e$ is exact. For (c), $\Phi^{\rm comb}\cong\Z/|S_s|$ via the degree, and $e,e^*$ are adjoint idempotents for the perfect pairing $\delta_{ee'}$.
\end{proof}
Thus the integral subtleties at finite level (Eisenstein cokernels, Mazur's $\mathrm{num}\frac{q-1}{12}$) are non-ordinary and disappear in Hida's limit.

Put $\cM^{\rm ord}_B(U):=\varprojlim_se\,\Z_p[X_B(UU_s)]$ (pushforward) for $U$ of level prime to $p$.
\begin{proposition}\label{prop:free}
(a) $e\Z_p[X_B(UU_s)]$ is free over $\Z_p[\Delta_s]$ \stproved{} (free action, $e$ a polynomial in $U_p$, $\Z_p[\Delta_s]$ local).
(b) $\cM^{\rm ord}_B(U)$ is $\Lambda$-free of finite rank with control $\cM^{\rm ord}_B\otimes_\Lambda\Z_p[\Delta_s]\cong e\Z_p[X_B(UU_s)]$ \stcited{} (Hida's argument; [LV], [Hs]).
\end{proposition}
(b) is not formal: fibres of $S_{s+1}\to S_s$ have size $p^2$ while $\ker(\Delta_{s+1}\to\Delta_s)$ has order $p$.

\begin{theorem}[$\Lambda$-adic definite Ribet sequence]\label{thm:lambdaribet}
Let $q\nmid Dp$ and $U$ of level prime to $pq$ with $U_q=\GL_2(\Z_q)$. Put $\cX^{\rm ord}_{Dq}(U):=\varprojlim_se\,H_1(\Gamma_{UU_s},\Z_p)$. Then
\[0\to\cX^{\rm ord}_{Dq}(U)\to\cM^{\rm ord}_B(U_0(q)U)\xrightarrow{\alpha}\cM^{\rm ord}_B(U)^{\oplus2}\to0\]
is an exact sequence of $\Lambda$-modules, equivariant for the Hecke operators away from $q$. Exactness \stproved{} (Thm.~\ref{thm:ribet}, Lemma~\ref{lem:ekills}, exactness of $e$ and of $\varprojlim$ on compact modules, given a $U_p$-normalisation commuting with pushforward); $\Lambda$-freeness of the three terms \stpmc{} (Prop.~\ref{prop:free}(b)); with $V_k$ after localisation at a non-Eisenstein $\m$ \stsketch.
\end{theorem}
\noindent Rank check \stproved: at $s\ge1$, $\rk H_1(\Gamma_{UU_s})=|E|-|V|+1=(q-1)|S_s|+1$.

\emph{Pairing.} The combinatorial pairing of Lemma~\ref{lem:ekills}(c), twisted along the tower (Hida/Ohta twist by $U_p^{-s}w_{p^s}$), gives a perfect $\Lambda$-pairing on $\cM^{\rm ord}_B$ \stsketch. For a primitive branch $F$ and non-Eisenstein $\m$ satisfying Hsieh's residual condition, $\cM^{\rm ord}_{B,\m}$ is free of rank one over the Gorenstein ring $\T^{\rm ord}_{B,\m}$, and the self-pairing of a generator of its $F$-eigenline generates the congruence ideal $\eta^B_F$ of $F$ in $\T^{\rm ord}_{B,\m}$ [Hs, Thm.~4.5, Def.~4.12, Rem.~7.8] \stcited. This is a statement about the definite algebra alone.

%======================================================================
\section{The comparison hypothesis}\label{sec:hypx}
%======================================================================

\subsection{Statement}
\begin{hypothesis}[$\mathrm{Hyp.X}(M;q)$]\label{hyp:X}
Data: $p\ge5$, $q\ne p$, $(M,pq)=1$, $B=B_q$; for $s\ge0$ put $\Gamma_s=\Gamma_1(Mp^s)\cap\Gamma_0(q)$, $X_s$ its modular curve, $J_s=\Jac X_s$, $S_s=X_{B_q}(U_1(Mp^s))$, and let $\mathcal H_s$ be the abstract Hecke algebra generated by $T_\ell$ ($\ell\nmid Mpq$), $U_\ell$ ($\ell\mid Mp$), $\langle d\rangle$ and $U_q$. There exist, for all $s$:
\begin{enumerate}[label=(X\arabic*)]
\item an isomorphism of $\mathcal H_s$-modules $\varphi_s:X_q(J_s)\xrightarrow{\sim}\Z[S_s]^0$ (Picard functoriality on the left, \S\ref{sec:notation} on the right, $U_q\mapsto\epsilon[\Pi_q^{\pm1}]$ with a universal sign $\epsilon$ and exponent to be fixed);
\item compatibility of pairings: $\varphi_s$ carries Grothendieck's monodromy pairing to $\langle e,e'\rangle=w_e\delta_{ee'}$;
\item naturality: for $s'\ge s$, $\varphi$ intertwines the map induced by $J_s\to J_{s'}$ with pushforward $\Z[S_{s'}]^0\to\Z[S_s]^0$;
\item (Galois part) a $\mathcal H_s$-stable filtration $T_t\subset T^{I_q}\subset T_pJ_s$, stable under a decomposition group at $q$, with graded pieces $\Hom(X_q(J_s),\Z_p(1))$, $T_pJ(\Gamma_1(Mp^s))^{\oplus2}$, $X_q(J_s)\otimes\Z_p$, monodromy given by (X2), transported to $H^1_{\rm par}(\Gamma_s,\Z_p)(1)$ by the comparison of Tate modules with Betti cohomology.
\end{enumerate}
$\mathrm{Hyp.X.char}:=$(X1)--(X3); $\mathrm{Hyp.X.gal}:=$(X4); $\mathrm{Hyp.X}:=$ all four. \sthyp
\end{hypothesis}
At $\Gamma_0$-type level (X1)--(X4) are the classical statements (G1)--(G4) of \S\ref{ss:geomroute} \stcited; at $\Gamma_1(Mp^s)$-level with $U_p$, diamonds and naturality in $s$ they are expected by the same proofs \stsketch. (X1)--(X3) can be phrased with Jacobians and character groups only; étale cohomology (or Tate modules with Galois action) is needed only for (X4).
In Lean the hypothesis is a single \texttt{def Hyp.X.char : Prop} (and one for the Galois part), taken as an explicit hypothesis of the theorems that use it, never as an axiom.

\begin{proposition}[$\Lambda$-adic version]\label{prop:Xlambda}
$\mathrm{Hyp.X.char}$, Lemma~\ref{lem:ekills} and Prop.~\ref{prop:free} give a $\Lambda[\T]$-isomorphism $\varprojlim_se(X_q(J_s)\otimes\Z_p)\cong\cM^{\rm ord}_B(M)$, carrying the $\Lambda$-adic monodromy pairing to the twisted Brandt pairing. With Hida's freeness of $\cT:=\varprojlim eT_pJ_s$ [Hi86a], $\mathrm{Hyp.X.gal}$ gives a $\Lambda$-adic monodromy filtration with graded pieces $\Hom_\Lambda(\cM^{\rm ord,*}_B,\Lambda)(1)$, $\cT(M)^{\oplus2}$, $\cM^{\rm ord,\dagger}_B$. Formal passage \stproved; twist \stsketch.
\end{proposition}
In particular, under $\mathrm{Hyp.X}$ there are $\Lambda$-adic Hecke-equivariant maps in both directions,
\[\Hom_\Lambda(\cM^{\rm ord,*}_B,\Lambda)(1)\hookrightarrow\cT\qquad\text{and}\qquad\cT\twoheadrightarrow\cM^{\rm ord,\dagger}_B,\]
whose composite is the (ordinary, perfect) monodromy map: an integral, non-semisimplified, cohomological analogue of Chenevier's $\JL_p$ on the ordinary locus for $D=q$ \stsketch. (Chenevier's statement about forms is not addressed.)

\subsection{What the tree route has without the hypothesis, and what is lost}
\begin{center}\small
\begin{longtable}{@{}p{5.6cm}p{4.2cm}p{4.6cm}@{}}
\toprule
statement & without $\mathrm{Hyp.X}$ & with $\mathrm{Hyp.X}$\\\midrule
$\Lambda$-adic Hecke algebras $\T^{D\text{-new}}\cong\T_B$ & \stcited{} [LV], [Hs] & same\\
$\Lambda$-adic character group, faithful action & definite side only & $D=q$ \stsketch\\
$\Lambda$-adic monodromy filtration; maps both ways & lost (Galois) & from (X4)\\
$e$ kills degree, vertex, $\Phi$; perfect pairing & \stproved{} (combinatorial) & $\Phi^{\rm comb}=\Phi_q(J)$ via (X2)\\
$\Lambda$-adic Ribet sequence & \stproved{} (definite form) & $\GL_2$ form needs (X) at a second prime and Čerednik--Drinfeld\\
$\eta^B_F$ as a self-pairing & \stcited{} (definite) & equals monodromy pairing via (X2)\\
$\Lambda$-adic defect formula (\S\ref{sec:defect}) & \stpmc{} for $D=q$ & not needed\\
\bottomrule
\end{longtable}
\end{center}
Without geometry one loses: (L1) the Galois action at $q$ (inertia filtration, $\mathrm{Frob}_q=qU_q$ on the torus), hence the maps above; (L2) a geometric proof of the reverse Tamagawa inclusion (it is recovered by level lowering, \S\ref{sec:defect}); (L3) the statement that a tree quotient \emph{is} a Shimura curve; (L4) $\Phi^{\rm comb}=\Phi_q(J)$; (L5) the dependence on $\iota$ (which was never Betti-canonical).

\subsection{A substitute at good maximal ideals}
Let $\m$ be non-Eisenstein, satisfying Hsieh's residual condition, $s\ge1$, ordinary parts. $\mathrm{Hyp.X}^{\rm iso}(s)$: $e\Z_p[S_s]_\m\cong eH^1_{\rm par}(\Gamma_s,\Z_p)^{\pm}_\m[\fa]$ as modules over the $q$-new Hecke algebra $\T'$ (the $q$-new \emph{sub}module), compatibly in $s$. Both sides are free of rank one over $\T'$ (multiplicity one for $eH^1_\m$ and Gorensteinness of $\T'$; [Hs, Thm.~4.5]), so $\mathrm{Hyp.X}^{\rm iso}$ holds, non-canonically \stsketch. It is a \emph{substitute} for (X1), not a consequence of JL: (X1) concerns the quotient lattice $X_q(J_s)\otimes\Z_p=T_pJ_s/T^{I_q}$. The two lattices differ exactly by the old/new congruence module:

\begin{proposition}[$q$-new sub- versus quotient lattice]\label{prop:subquot}
Let $H$ be free of rank $r$ over $\T$, with $\T$ $\Lambda$-Gorenstein and $\T''$ $\Lambda$-free (notation of \S\ref{sec:defect}). Then $H[\fa]=\fb H\cong\fb^r$, $H/\fa H\cong\T'^{\,r}$, and $H[\fa]\to H/\fa H$ is injective with cokernel $(\T'/\cc_{\T'})^r$, $\cc_{\T'}=(\fa+\fb)/\fa$. \stproved
\end{proposition}
So $\mathrm{Hyp.X}^{\rm iso}$ suffices for congruence- and Fitting-ideal statements at good $\m$, but not for Eisenstein or non-Gorenstein $\m$, for pairings beyond $\eta^B$, or for anything Galois. A Betti-canonical \emph{map} replacing $\mathrm{Hyp.X}$ is not expected: after inverting $q$, the cuspidal cohomology of $\GL_2$ sits in degree $2$ and that of the definite group in degree $1$ (or $0$ when $q\mid D$), and the gap never closes \stsketch.

%======================================================================
\section{The integral JL defect}\label{sec:defect}
%======================================================================

\subsection{Setting}
Standing hypotheses:
\begin{itemize}
\item[(H0)] $p\ge5$, $q\ne p$, $(M,pq)=1$, $N:=qM$, $p\nmid\varphi(N)$ (so $q\not\equiv1\bmod p$).
\item[(H1)] $\T$ is Hida's full ordinary cuspidal Hecke algebra of tame level $\Gamma_0(N)$ over $\Lambda$ (all $T(n)$, so it contains $U_q$), in a weight component $k_0\bmod(p-1)$, localised at a non-Eisenstein maximal ideal $\m$. $F$ is a primitive ordinary $q$-new branch, $\lambda_F:\T\to I_F$, with $I_F$ a normal domain finite over $\Lambda$ (hence $\Lambda$-free).
\item[(H2)] $\bar\rho=\bar\rho_\m$ is absolutely irreducible on $G_{\Q(\sqrt{p^*})}$ and $p$-distinguished.
\item[(Im)] $p$ divides $\#\bar\rho(G_\Q)$.
\item[(P)] Either $k_0\not\equiv2\bmod(p-1)$ (then $F$ has an arithmetic point of weight $k_P\in[3,p-1]$, a \emph{Kim--Ota point}), or $k_0\equiv2$ and $F$ has a weight-two point which is the ordinary $p$-stabilisation of a newform of level $N$.
\end{itemize}
Let $\T'$, $\T''$ be the images of $\T$ on $q$-new and $q$-old ordinary $\Lambda$-adic cusp forms, $\fa=\ker(\T\to\T')$, $\fb=\ker(\T\to\T'')$; $R=\T\otimes_\Lambda I_F$, $R'$, $R''$ similarly; $\lambda:R\to I_F$, $\lambda':R'\to I_F$ induced by $\lambda_F$. Put
\[\eta_F:=\eta(R,\lambda),\qquad\eta^B_F:=\eta(R',\lambda'),\qquad h_F:=\eta_F/\eta^B_F.\]
Via $\T'\cong\T_B$ ($\Lambda$-adic JL of Hecke algebras), $\eta^B_F$ is the congruence ideal of $F$ on the definite side; $h_F$ measures the failure of an integral JL isomorphism of lattices. This is the defect that [IP] predicts must exist.

\begin{definition}[combinatorial Tamagawa ideal]\label{def:ctam}
$\cc:=\mathrm{im}(\fb R\to R')$, so $R'/\cc\cong R/(\fa R+\fb R)$ is the $\Lambda$-adic old/new congruence module at $q$; $\cc_q(F):=\lambda'(\cc)$.
\end{definition}
The module $R/(\fa R+\fb R)$ is Diamond's $\Lambda$-adic old/new congruence module [Dia89], who computes it on the \emph{old} branch. On the $\GL_2$ side its fixed-weight version is computed by $\coker(\Psi^*\Psi)$ of \S\ref{ss:amalgam}; on the definite side the analogue for $B_r$ with tree at $q$ is $\coker(\alpha\alpha^*)$ of Thm.~\ref{thm:ribet} \stsketch.

\begin{definition}[Tamagawa ideal]
$\tq_q(F)\subseteq I_F$ is the ideal generated by the entries of $\rho_F(\sigma)-1$, $\sigma\in I_q$, for the Galois representation $\rho_F$ on $\cT_F:=H\otimes_{\T,\lambda_F}I_F$, $H$ the ordinary $\Lambda$-adic étale cohomology localised at $\m$ (free of rank two over $\T$ under (H2) [Wi95], [Hs]). $\rho_F|_{I_q}$ is unipotent and factors through $\Z_p(1)$, so $\tq_q(F)$ is generated by the entries of $\rho_F(\tau)-1$ for a generator $\tau$ \stproved. At a $q$-new arithmetic point $P$, $P(\tq_q(F))=(\varpi^{t_{F_P}(q)})$ with $t_f(q)$ the largest $t$ such that $\rho_f\bmod\varpi^t$ is unramified at $q$ (Kim--Ota's Tamagawa exponent [KO]) \stproved.
\end{definition}

\subsection{Commutative algebra}
\begin{lemma}[base change of congruence ideals]\label{lem:S}
Let $S$ be a finite free $A$-algebra with $\Hom_A(S,A)$ free of rank one over $S$ ($S$ is $A$-Gorenstein), $\mu:S\to A$ an $A$-algebra map, $\varphi:A\to A'$. Then $\Ann_S\ker\mu=A\theta$ with $\ell(\theta x)=\mu(x)$ for a generator $\ell$, $\eta(S,\mu)=(\mu(\theta))$ is principal, and $\eta(S\otimes A',\mu\otimes A')=\varphi(\eta(S,\mu))A'$. \stproved
\end{lemma}
\begin{lemma}[Gorenstein from one fibre]\label{lem:G}
If $A$ is Noetherian local, $S$ local and finite free over $A$, $\mathfrak p\subseteq\m_A$, and $S/\mathfrak pS$ is $A/\mathfrak p$-Gorenstein, then $S$ is $A$-Gorenstein. \stproved{} (Nakayama.)
\end{lemma}
\begin{lemma}[new/old duality]\label{lem:D}
Let $A$ be a Noetherian local domain, $S$ local, finite free and $A$-Gorenstein, $\fa,\fb\subseteq S$ with $S/\fa$, $S/\fb$ torsion-free, $\fa\cap\fb=0$ and $S_K=(S/\fa)_K\times(S/\fb)_K$. Then $\Ann(\fa)=\fb$, $\Ann(\fb)=\fa$; if $S/\fb$ is $A$-free then $S/\fa\cong\fb^\vee$ is $A$-free, and these statements commute with base change. \stproved
\end{lemma}

\begin{theorem}[the combinatorial half]\label{thm:E}
In the setting of Lemma~\ref{lem:D} with $S=R$, $S'=R'$, let $\lambda$ factor through $\lambda'$. If $R'$ is $I_F$-Gorenstein, then
\[\eta_F=\cc_q(F)\cdot\eta^B_F,\qquad\text{i.e. }(h_F)=\cc_q(F).\]
Without the Gorenstein property of $R'$ the equality holds after localisation at every height-one prime of $I_F$; and in general $\cc\cong\Hom_{I_F}(R',I_F)$, so $\cc$ is principal iff $R'$ is Gorenstein. \stproved
\end{theorem}
\begin{proof}[Proof sketch]
Let $\theta\in R$ generate $\Ann_R\ker\lambda$ (Lemma~\ref{lem:S}). Since $\fa\subseteq\ker\lambda$, $\theta\in\Ann(\fa)=\fb$; its image $x\in R'$ kills $\ker\lambda'$. The isomorphism $\fb\cong R'^\vee$, $\beta\mapsto\ell(\beta\,\cdot)$, gives $\lambda'(\cc)=\{\psi(x):\psi\in R'^\vee\}$. If $R'$ is Gorenstein, $\Ann_{R'}\ker\lambda'=I_F\theta'$ is a direct summand, $x=h\theta'$, and the content of $x$ is $(h)$.
\end{proof}

\begin{proposition}[Galois comparison]\label{prop:galois}
Under (H1) and absolute irreducibility of $\bar\rho$, with $\cT_F$ free of rank two: $\tq_q(F)\subseteq\cc_q(F)$. \stpmc{} (Galois representation of tame level $M$, unramified at $q$; Chebotarev).
\end{proposition}
\begin{proof}[Proof sketch]
Put $A=I_F/\cc_q(F)$. Modulo $\cc$, $\lambda$ factors through the $q$-old quotient, whose $T_\ell$ come from tame level $M$; so $\tr\rho_F\otimes A$ is the trace of a representation unramified at $q$, i.e.\ $\tr\bar\rho_A(g\sigma)=\tr\bar\rho_A(g)$ for $\sigma\in I_q$. Since $A[G_\Q]\to M_2(A)$ is surjective (Burnside and Nakayama) and the trace form on $M_2(A)$ is perfect, $\bar\rho_A(\sigma)=1$. No pseudo-representation theorem is needed.
\end{proof}

\subsection{The defect formula for $D=q$}
\begin{theorem}[$\Lambda$-adic defect formula]\label{thm:defect}
Assume (H0), (H1), (H2), (Im) and (P). Then
\[(h_F)\;=\;\cc_q(F)\;=\;\tq_q(F)\qquad\text{as principal ideals of }I_F.\]
\stpmc. Cited inputs: Hida's freeness and control [Hi86b, Thm.~3.1, Cor.~3.2, Cor.~3.3]; Hsieh's theorem that $\T$ and the $q$-new quotient $\T'\cong\T^{B}$ are $\Lambda$-Gorenstein [Hs, \S3.3, Thm.~4.5, Rem.~7.8], resting on [Wi95], [MW], [Hel]; freeness of ordinary cohomology [Hs], [Wi95]; and, at the point of (P), level lowering modulo $\varpi^t$ (below), resting on an $R=\T$ theorem for new-at-$\Sigma$ deformation rings [Dim09, Thm.~6.6], [DFG], [KO, \S\S2--3] and Böckle's presentation.
\end{theorem}
\begin{proof}[Proof outline]
(1) $\T$ and $R'$ are Gorenstein (Hsieh; $R'=\T'\otimes I_F$, Lemma~\ref{lem:S}), so $(h_F)=\cc_q(F)$ by Thm.~\ref{thm:E}.
(2) $\tq_q(F)\subseteq\cc_q(F)$ by Prop.~\ref{prop:galois}; write $\tq_q(F)=h_F\mathfrak d$.
(3) At the point $P$ of (P), with $f$ the newform of level $N$: $\T'$ is $\Lambda$-free with $\T'/P\T'$ the $q$-new quotient at weight $k_P$ (Lemma~\ref{lem:D} against the $\Lambda$-free old quotient $\T''\cong\T(M)[X]/(X^2-T_qX+c_q)$, localised), so $P(\cc_q(F))\cO\subseteq\cc_q(f)$, the fixed-weight old/new congruence ideal.
(4) \emph{Level lowering modulo $\varpi^t$}: $\cc_q(f)\subseteq(\varpi^{t_f(q)})$ for ordinary $f$ of weight $2\le k\le p-1$, any tame level $M$ with $p\nmid\varphi(qM)$, including $q\equiv-1\bmod p$, under (H2), (Im) \stpmc. It is proved by passing to the quotient that is new at the primes $\ell\,\|\,M$ where $\bar\rho$ is unramified and applying the $R=\T$ theorem there.
(5) Hence $P(\mathfrak d)$ is a unit, so $\mathfrak d=I_F$ ($I_F$ local), and the three ideals coincide.
\end{proof}

\begin{corollary}[all weights of the branch]\label{cor:allweights}
On a branch as in Thm.~\ref{thm:defect}, at every arithmetic point of weight $k'\ge3$ with trivial character (including $k'\ge p$), the fixed-weight old/new congruence ideal equals $(\varpi^{t_{f'}(q)})$ and $\ord_\varpi\eta_{f'}(N)-\ord_\varpi\eta_{f'}(M,q)=t_{f'}(q)$ (ring-theoretic congruence ideals of the full and the $q$-new Hecke algebras at weight $k'$). \stpmc{} (base change, Lemmas~\ref{lem:S}, \ref{lem:D}).
\end{corollary}
At fixed weight $2\le k\le p-1$ the identity $\ord\eta_f(N)-\ord\eta_f(M,q)=t_f(q)$ is due to Kim--Ota [KO] (Pollack--Weston [PW] at weight two). What the tree route adds is the $\Lambda$-adic statement on the $q$-\emph{new} branch, the identification of the defect with the combinatorially computable $\cc_q$, and the extension to all weights of the branch.

\begin{remark}[open cases]
Branches with only $p$-new weight-two points ((Comb) holds; the Galois half needs $R=\T$ at level $Mp$); trivial primes ($q\equiv1\bmod p$ with $\bar\rho(\mathrm{Frob}_q)$ scalar), where $\T'$ fails to be Gorenstein [BKM] and $\cc$ is not principal (Thm.~\ref{thm:E}); dropping (Im), which requires [DFG]; several primes in $D$ (\S\ref{sec:general}).
\end{remark}

%======================================================================
\section{Several ramified primes}\label{sec:general}
%======================================================================

\begin{proposition}[no tree-only chain of module maps]\label{prop:noChain}
Let $\mathcal C$ be the class of Hecke-equivariant module maps generated by the identifications of Prop.~\ref{prop:quot} for a fixed definite algebra, the maps of Thm.~\ref{thm:ribet}, compositions, kernels and cokernels. (a) For definite $B'$ and $q\nmid\mathrm{disc}B'$ the cells of tree quotients are class sets of $B'$; for $q\mid\mathrm{disc}B'$ there is no tree ($\cO'_q[\frac1q]^1=\cO'^1_q$ is compact). (b) For $\omega(D)\ge3$, no sequence of operations in $\mathcal C$ applied to definite algebras of smaller discriminant produces the Brandt module of $B_D$. \stproved{} (the operations never change the algebra).
\end{proposition}
So for $\omega(D)\ge3$ each change of definite algebra needs either a Shimura-curve comparison (a version of $\mathrm{Hyp.X}$ for Shimura curves, i.e.\ Deligne--Rapoport-type models of Shimura curves) or JL as an input. The geometric chain for $D=q_1q_2q_3$ is:
\begin{center}\small
\begin{tabular}{@{}p{5.6cm}p{3.6cm}p{4.6cm}@{}}
\toprule
link & tree-only & geometric input\\\midrule
$X_{q_1}(J_0(MD))\cong\Z[X_{B_{q_1}}(U_0(q_2q_3M))]^0$ & definition & $\mathrm{Hyp.X}$ at $q_1$\\
Ribet collapse for $B_{q_1}$, tree at $q_2$ & Thm.~\ref{thm:ribet} & none\\
tree $H_1$ $=X_{q_2}(J^{q_1q_2}(q_3M))$ & lost & Čerednik--Drinfeld\\
$X_{q_2}(J^{q_1q_2})\leftrightarrow X_{q_3}(J^{q_1q_2})$ & no substitute & only through $T_\ell J$ (Galois)\\
$X_{q_3}(J^{q_1q_2}(q_3M))\cong\Z[X_{B_D}(M)]^0$ & definition & Shimura-curve version of $\mathrm{Hyp.X}$\\
\bottomrule
\end{tabular}
\end{center}
This concerns maps of modules. At the level of Hecke algebras and ideals the $\Lambda$-adic isomorphism $\T^{D\text{-new}}\cong\T_{B_D}$ holds for all $D$, and a defect formula $(h_F)=\prod_{q\mid D}\tq_q(F)=\prod_i\cc^{(i)}_{q_i}(F)$, via a filtration by $q_1\cdots q_i$-new quotients, is \stsketch{}/\stconj; whether the intermediate quotients need to be Gorenstein is open.

%======================================================================
\section{Simplicial shape of the route}\label{sec:simplicial}
%======================================================================
All objects of \S\S\ref{sec:tree}--\ref{sec:lambda} are finite simplicial orbifolds or their cohomology:
\begin{itemize}
\item \emph{Definite level $U$}: the $0$-dimensional orbifold $\coprod_{[x]\in X_B(U)}B\Gamma_x$; its cohomology with coefficients $V$ is $M(U;V)$.
\item \emph{Definite $S$-arithmetic stage}: the graph of finite groups $\Gamma_U\backslash\backslash\cT_q$, with vertex groups at two copies of $X_B(U)$, edge groups at $X_B(U_0(q)U)$, and the degeneracies $s_0,s_1$ as structure maps. Since $\cT_q$ is contractible, it is a $1$-dimensional orbifold model of $B\Gamma_U$ (a Moerdijk--Pronk-type simplicial orbifold after barycentric subdivision), and its cohomology with local coefficients is $H^*(\Gamma_U,V)$ (Prop.~\ref{prop:coh}). At $U_1(p^s)$-level ($s\ge1$) all stabilisers are trivial, $\Gamma_U$ is free, and the quotient is an honest graph.
\item \emph{$\GL_2$ arithmetic stage}: $\Gamma_0(N)\backslash\Hb$, realised simplicially (e.g.\ by a finite triangulation of a fundamental domain of $\SL_2(\Z)$ with Shapiro induction, as in HACS).
\item \emph{$\GL_2$ $S$-arithmetic stage}: the amalgam of \S\ref{ss:amalgam}, a graph of spaces over the same two-vertex graph shape $[\,\bullet\ \bullet\leftleftarrows\text{edge}\,]$ as the definite stage.
\end{itemize}
The two $S$-arithmetic stages have the same combinatorial shape (two vertex objects, one edge object, two degeneracies), with definite class sets in place of modular curves. The zig-zag
\[\underbrace{X_B(U)\rightleftarrows X_B(U_0(q)U)}_{\text{0-dim}}\ \leadsto\ \underbrace{\Gamma_U\backslash\backslash\cT_q}_{\text{1-dim}}\ \overset{\mathrm{Hyp.X}}{\longleftrightarrow}\ \underbrace{J(\Gamma_0(q)\cap\Gamma_1(M))\text{ at }q}_{\GL_2}\]
consists of explicit simplicial maps in its first part; the comparison in its last part is $\mathrm{Hyp.X}$, i.e.\ the Deligne--Rapoport graph of the modular curve, which is the DR graph $\G^{\rm DR}_q$ of Def.~\ref{def:DR}. Making the last arrow a map of explicit simplicial objects (rather than a hypothesis on character groups) is the main open structural question of the route.

%======================================================================
\section{HACS: computations and test cases}\label{sec:hacs}
%======================================================================

\subsection{What HACS is}
HACS [HACS] is a C++ program (with optional CUDA) for the cohomology of $\SL_2(\Z)$, its congruence subgroups and $\GL_2$, with coefficients $\Sym^n$ of the standard representation over $\F_p$, $\Z/\ell^k$ and $\Z$, together with the Hecke action, in exact modular arithmetic. Its features relevant here:
\begin{itemize}
\item an explicit simplicial model of $\SL_2(\Z)\backslash\Hb$ (a $7$-vertex triangulation of a truncated fundamental domain with gluing and stabiliser relations), and a Moerdijk--Pronk simplicial set computing group cohomology in all degrees and over all rings, with a Borel--Serre-type boundary and the long exact sequence of the pair;
\item level $\Gamma_0(N)$, $\Gamma_1(N)$ and nebentypus by Shapiro induction on that model;
\item Hecke operators $T_\ell$, $U_\ell$ in all degrees (chain-level and Manin--Merel--Heilbronn), the ordinary projector and Hida theory at small levels;
\item integral structure: torsion, integral Hecke algebras, saturation, congruence modules $C_\T(A,B)=\T/(J_A+J_B)$ with certified bounds, recorded in signed data records.
\end{itemize}
HACS has no quaternion algebras, no $S$-arithmetic groups beyond what can be assembled from $\SL_2(\Z)$-data, and no JL.

\subsection{Computations already carried out}\label{ss:done}
All with $V=\Sym^{k-2}$, classical Hecke normalisation. Here $q$ denotes the level prime (written $D$ in the first item, where $D=q$).
\begin{enumerate}[label=(C\arabic*)]
\item \emph{Old/new congruence modules at level $q$, weight $12$} \stcomp. For $\Gamma_0(3)$, $k=12$: $H^1(\Gamma_0(3);\Sym^{10}\Z^2)=\Z^8\oplus\Z/2\oplus\Z/4\oplus(\Z/3)^4$; the congruence module between the $3$-old ($\Delta$) and $3$-new forms of the full Hecke algebra is $\Z/17$ at $17$ (certified), matching $\tau(2)-a_2(g)=-102=-2\cdot3\cdot17$ for the $3$-new form $g$. For $q=5$: primes $29$ and $131$, both with exponent $1$. For $q\in\{3,5,7,11,13\}$, $k=12$, the primes $l>12$ of the anemic old/new congruence module are $\{17\}$, $\{29,131\}$, $\{1051\}$, $\{233,34267\}$, $\{31,191,5323\}$, exactly the primes $l$ with $l\mid a_q(\Delta)^2-q^{10}(q+1)^2$ (level raising). A table of predicted level-raising primes for $q\le13$, $k\le26$ is available.
\item \emph{Ordinary ranks} at level $3$ for $k\le62$, $p\in\{5,7,11,13\}$, constant on weight classes mod $p-1$ \stcomp, and a mod-$17$ persistence test along $k=12+16j$ (certified ``no'' at $k=16,\dots,22$, not on the $17$-adic family) \stcomp.
\item \emph{The amalgam $\SL_2(\Z[1/q])=\SL_2(\Z)*_{\Gamma_0(q)}\SL_2(\Z)$}, Mayer--Vietoris with $\Sym^{k-2}$ over $\Z_{(l)}$, $l>k$, $l\ne q$, for $q\in\{3,5,7,11,13\}$, $k=4,\dots,24$ \stcomp:
\begin{itemize}
\item $H^0=H^1(\SL_2(\Z[1/q]),V)=0$, $H^{\ge3}=0$, and $H^2=\coker\Psi$ is free of rank $2\dim S_k(\Gamma_0(q))^{q\text{-new}}$ plus torsion;
\item the torsion of $H^2$ is $0$ or a single $\Z/l$ with $l=k+1$, Eisenstein, in every cell; there is no torsion at the cuspidal congruence primes ($17$ for $q=3$; $29$, $131$ for $q=5$, $k=12$): consistent with an Ihara lemma for $\GL_2$;
\item with the transfer $\Psi^*$, $\Psi^*\Psi=\left(\begin{smallmatrix}q+1&-T_q\\-q^{2-k}T_q&q+1\end{smallmatrix}\right)$ on $H^1(\SL_2(\Z),V)^2$, and $\coker(\Psi^*\Psi)$ at $\m_\Delta$ is $(\Z/l^v)^2$ (one copy per sign) with $v=v_l(\tau(q)^2-(q+1)^2q^{10})$: $(\Z/17)^2$ at $q=3$, $(\Z/29)^2$ and $(\Z/131)^2$ at $q=5$; the Eisenstein part has length $v_l((1-q^{k-2})(1-q^k))$;
\item $T_\ell$ ($\ell\ne q$, $\ell\le13$) on all terms; independent cross-checks against HACS: $2447$ comparisons, no mismatch.
\end{itemize}
This is the $\GL_2$ end of the $S$-arithmetic stage: $\coker(\Psi^*\Psi)_\m$ is the Betti incarnation of the old/new congruence module $\cc_q$.
\item \emph{Independent computations (PARI/GP, not HACS)} \stcomp. (i) For $q\in\{2,3,5,7,11,13\}$, $qM\le70$, even $4\le k\le26$ (resp.\ $30$), all $p\le10^6$: $\eta_F$, $\eta^B$ (ring-theoretic, on the $q$-new quotient) and $\cc_q$ at maximal ideals with at least two $q$-new eigenforms. In all $15$ cases with old support inside the range $p>k$, $p\nmid\varphi(N)$, $q\not\equiv\pm1$, non-Eisenstein, $v(h)=v(\cc_q)$; failures occur only at Eisenstein or twisted-Eisenstein (reducible $\bar\rho$) maximal ideals, consistent with $(h)\supseteq\cc_q$. By Thm.~\ref{thm:E} this tests the Gorenstein property of the full Hecke algebra at fixed weight, not the Galois side. (ii) Weight-two Brandt matrices from supersingular isogeny graphs for $q\in\{11,17,19,23,37,53\}$ agree with $T_\ell$ on $S_2(\Gamma_0(q))$ and satisfy the mass formula.
\end{enumerate}

\subsection{What the route hopes to add to HACS}
\begin{itemize}
\item a \emph{definite module}: class sets of maximal and Eichler orders of $B_D$ with stabilisers, Brandt modules with $\Sym^n$ coefficients, Hecke operators including $U_p$ at Iwahori level and $W_\ell$ ($\ell\mid D$), the Brandt (Gram) pairing and its congruence modules $\eta^B$; validated by mass formulas and JL trace comparisons against HACS level-$D$ new spaces;
\item the \emph{definite $S$-arithmetic stage}: the tree quotient $\Gamma_U\backslash\backslash\cT_q$ as incidence matrices $s_0,s_1$, its Mayer--Vietoris (Ribet/Ihara) sequence, $\coker(\alpha\alpha^*)$; this is the same shape of computation as (C3), with finite vertex groups;
\item $p$-power towers on both sides with $U_p$ at level $qp^s$ (missing in HACS on the $\GL_2$ side), the ordinary projector and $\Lambda$-adic control;
\item comparison data for $\mathrm{Hyp.X}$ beyond isomorphism classes: Gram matrices versus $\GL_2$ pairings at Eisenstein and non-Gorenstein maximal ideals.
\end{itemize}
The existing integral machinery of HACS (sparse elimination mod $\ell^K$, saturation, congruence-module records) applies unchanged to these finite free modules with commuting integral operators.

\subsection{Test cases}\label{ss:tests}
Validation items check code against a theorem; ``new'' items have no theorem predicting the result.
\begin{enumerate}[label=\textbf{T\arabic*}]
\item \textbf{[validation] DR graph and $\Phi^{\rm comb}$, level $1$.} $q\in\{11,23,37,101,131\}$: compute $X_{B_q}(1)$, the weights, and the Smith form of the weighted pairing on $\Z[E]^0$. Prediction: $\Phi^{\rm comb}\cong\Z/\mathrm{num}\frac{q-1}{12}$ ($5,11,3,25,65$).
\item \textbf{[validation] Ribet collapse, trivial coefficients.} (a) $B_2$, $q=7$: $\rk H_1=1=g(X^{14})$. (b) $B_2$, $q=13$: $\rk H_1=2=g(X^{26})=\dim S_2(\Gamma_0(26))^{\rm new}$. (c) $B_{11}$, $q=2$: two vertices per type, $\rk H_1=0=g(X^{22})$, so $\alpha$ is an isomorphism over $\Z$; compute $\coker(\alpha\alpha^*)$ with both the weighted and the unweighted adjoint. Expectation: order $|a_2(11a)^2-9|=5$ up to factors $2,3$.
\item \textbf{[validation] $\Lambda$-adic Ribet at deep level.} $B_3$, tree at $q=2$, $p=5$, $U_s=U_1(5^s)$, $s=1,2,3$: $|S_s|=2\cdot25^{s-1}$, $|E_s|=3|S_s|$, $\rk H_1=|S_s|+1$ ($s=1$: $3$, matching the genus of the corresponding Shimura curve of discriminant $6$ at $\Gamma_1(5)$-level). Apply $e$ mod $5^3$; predict exactness, $\Z/5^3[\Delta_s]$-freeness of $s$-independent rank, degree $\equiv0$ on the ordinary part.
\item \textbf{[validation/prediction] $\cc_q$ along a branch.} $q=3$, $p=17$, the branch through $g\in S_{12}^{\rm new}(\Gamma_0(3))$, $g\equiv\Delta\bmod17$: compute $\coker(\Psi^*\Psi)_\m$ at $k=12,28,44,60,124$. Prediction: $\ord_{17}\cc_3=1,1,1,1,2$. The $q$-new quotient has rank one here, so $(h)=\cc$ is automatic; this tests the Betti computation of $\cc_q$ and, where $t_{G_k}(3)$ is computed independently, Cor.~\ref{cor:allweights} for $k\ge p$.
\item \textbf{[new] Combinatorial half with $\rk_\Lambda\T'_\m\ge2$.} Level-raising triples $(q,k,p)$ with $3\le k\le p-1$, $p\nmid\varphi(q)$, $q\not\equiv\pm1\bmod p$, e.g.\ $(13,18,19)$, $(5,18,23)$, $(5,16,31)$, $(13,12,31)$, $(3,18,53)$, $(3,18,61)$, $(13,16,101)$: compute the mod-$p$ dimension of the $q$-new Brandt module of $B_q$ with $\Sym^{k-2}$ at $\m$, and where it is $\ge2$ compare $\eta_f(q)/\eta_f(1,q)$ with $\cc_q(f)$ \emph{on the definite side} (Gram matrix of the Brandt module). Repeat at $k+(p-1)$.
\item \textbf{[new] Higher-weight cokernel.} $H_0(\Gamma_U,V_k)=V_{k,U^1_p}$ for $B_3$, tree at $2$, $p\in\{5,7\}$, $k\le2p$: order and Hecke action; expectation: $T_\ell$ acts by $(1+\ell)\chi(\ell)^{\pm1}$ on $\det$-eigenpieces (Eisenstein). Pins the sign convention of Cor.~\ref{cor:ihara}(iii).
\item \textbf{[later] $\mathrm{Hyp.X}$ beyond isomorphism class.} With $U_p$ at level $qp$ on the $\GL_2$ side: compare the Brandt Gram matrix on $e\Z_p[S_s]^0_\m$ with the (Ohta-twisted) $\GL_2$ pairing on $eH^1_{\rm par}(\Gamma_s)^{\pm}_\m$ at Eisenstein or non-Gorenstein $\m$.
\end{enumerate}
T1--T3 and T6 are small; T4 needs the amalgam at weights up to $124$; T5 needs the definite module with $\Sym^n$.

%======================================================================
\section{Formalisation}\label{sec:lean}
%======================================================================
Target: Lean~4 with Mathlib, no \texttt{axiom}; deep inputs enter as explicit hypotheses \texttt{def Hyp.X : Prop}, registered and documented.
\begin{itemize}
\item \emph{Lean-friendly now}: Lemmas~\ref{lem:tree}, \ref{lem:stab}; Prop.~\ref{prop:phi} (finite linear algebra); Thm.~\ref{thm:ribet}(a) as a statement about two-vertex collapses of finite multigraphs (Mathlib's simple graphs are too weak; use quivers or incidence matrices); Cor.~\ref{cor:ihara}(ii); Lemma~\ref{lem:ekills}; the commutative algebra of \S\ref{sec:defect} (Lemmas~\ref{lem:S}--\ref{lem:D}, Thm.~\ref{thm:E}, Prop.~\ref{prop:subquot}) and the final squeeze of Thm.~\ref{thm:defect}.
\item \emph{One deep input for the combinatorial core}: strong approximation for $B^1$ at a split place, in the form $y=\beta gu$ (needed for Prop.~\ref{prop:quot}); existing Lean developments contain it in this shape. The tree itself is formalised [LM], including the Laplacian surjectivity on the full tree; the $\Gamma$-equivariant version with coefficients is not.
\item \emph{Avoidable input}: the equivariant spectral sequence (Prop.~\ref{prop:coh}) is not needed at $U_1(p^s)$-level, where $\Gamma_U$ is free and acts freely, so $H_1(\Gamma_U)=H_1$ of the quotient graph.
\item \emph{Mathlib gaps}: quaternion orders, reduced norm, class sets, mass formula; Brandt modules with Hecke operators; graphs of groups; Hida control; congruence ideals (elementary).
\item \emph{Registered hypotheses}: $\mathrm{Hyp.X.char}$ and $\mathrm{Hyp.X.gal}$ (\S\ref{sec:hypx}); Hida's control theorems; Hsieh's Gorenstein theorem; Galois representations of Hida families with their freeness; the $R=\T$ input of level lowering. None of these is geometric except $\mathrm{Hyp.X}$, and $\mathrm{Hyp.X}$ is not used for the defect formula.
\end{itemize}
By contrast (\S\ref{ss:cost}), the geometric route would need Deligne--Rapoport models, Néron models, Grothendieck's monodromy and Čerednik--Drinfeld uniformisation as hypotheses at every level, or a very large imported library.

%======================================================================
\section{Open points}\label{sec:open}
%======================================================================
\begin{enumerate}[label=O\arabic*.]
\item $\mathrm{Hyp.X}$ at $\Gamma_1(Mp^s)$-level with $U_p$, diamonds and naturality; the sign $\epsilon$ and exponent of $U_q\mapsto\epsilon\Pi_q^{\pm1}$; the $U_p$ normalisation; the Ohta twist of the pairing.
\item A map of explicit simplicial objects realising the last arrow of the zig-zag of \S\ref{sec:simplicial} (or a proof that only $\mathrm{Hyp.X}$-type statements are possible).
\item Higher weight: Hecke and $U_p$ action on $H_0(\Gamma_U,V_k)$; whether $e$ kills it.
\item Several primes in $D$: the defect formula as ideals (\S\ref{sec:general}), and Shimura-curve versions of $\mathrm{Hyp.X}$ if module maps are wanted.
\item The open cases of the defect formula (branches without (P), trivial primes, without (Im)).
\item Eisenstein and non-Gorenstein maximal ideals, where $\mathrm{Hyp.X}$ carries information that $\mathrm{Hyp.X}^{\rm iso}$ does not (Example~\ref{ex:mazur}, T1, T7).
\item The relation of the $\Lambda$-adic maps of \S\ref{sec:hypx} with Chenevier's eigencurve immersion away from the ordinary locus.
\end{enumerate}

%======================================================================
\begin{thebibliography}{HACS}\small
\bibitem[BC]{BC} J.-F.~Boutot, H.~Carayol, Uniformisation $p$-adique des courbes de Shimura: les théorèmes de Čerednik et de Drinfeld, Astérisque 196--197 (1991), 45--158.
\bibitem[BKM]{BKM} G.~Böckle, C.~Khare, J.~Manning, Wiles defect for Hecke algebras that are not complete intersections, Compositio Math.\ 157 (2021).
\bibitem[Br]{Br} K.~S.~Brown, \emph{Cohomology of groups}, GTM 87, Springer.
\bibitem[Če]{Ce} I.~V.~Čerednik, Uniformization of algebraic curves by discrete arithmetic subgroups of $\PGL_2(k_w)$ with compact quotient spaces, Mat.\ Sb.\ 100 (1976).
\bibitem[Ch]{Ch} G.~Chenevier, Une correspondance de Jacquet--Langlands $p$-adique, Duke Math.\ J.\ 126 (2005), 161--194.
\bibitem[DFG]{DFG} F.~Diamond, M.~Flach, L.~Guo, The Tamagawa number conjecture of adjoint motives of modular forms, Ann.\ Sci.\ ÉNS 37 (2004).
\bibitem[Dia89]{Dia89} F.~Diamond, On congruence modules associated to $\Lambda$-adic forms, Compositio Math.\ 71 (1989), 49--83.
\bibitem[Dim09]{Dim09} M.~Dimitrov, On Ihara's lemma for Hilbert modular varieties, Compositio Math.\ 145 (2009).
\bibitem[DR]{DR} P.~Deligne, M.~Rapoport, Les schémas de modules de courbes elliptiques, LNM 349, Springer 1973.
\bibitem[Dr]{Dr} V.~G.~Drinfeld, Coverings of $p$-adic symmetric regions, Funct.\ Anal.\ Appl.\ 10 (1976).
\bibitem[DT]{DT} F.~Diamond, R.~Taylor, Non-optimal levels of mod $l$ modular representations, Invent.\ Math.\ 115 (1994).
\bibitem[HACS]{HACS} HACS: cohomology of $\SL_2(\Z)$, $\GL_2$ and congruence subgroups with Hecke action, software, \url{https://github.com/drjanosch42/hacs}, Zenodo DOI 10.5281/zenodo.24118841; data: DOI 10.5281/zenodo.24118877.
\bibitem[Hel]{Hel} D.~Helm, On maps between modular Jacobians and Jacobians of Shimura curves, Israel J.\ Math.\ 160 (2007).
\bibitem[Hi86a]{Hi86a} H.~Hida, Galois representations into $\GL_2(\Z_p[[X]])$ attached to ordinary cusp forms, Invent.\ Math.\ 85 (1986).
\bibitem[Hi86b]{Hi86b} H.~Hida, Iwasawa modules attached to congruences of cusp forms, Ann.\ Sci.\ ÉNS 19 (1986), 231--273.
\bibitem[Hs]{Hs} M.-L.~Hsieh, Hida families and $p$-adic triple product $L$-functions, Amer.\ J.\ Math.\ (2021); arXiv:1705.02717.
\bibitem[Ill]{Ill} L.~Illusie, Réalisation $\ell$-adique de l'accouplement de monodromie d'après A.~Grothendieck, Astérisque 196--197 (1991), 27--44.
\bibitem[IP]{IP} A.~Ichino, K.~Prasanna, Periods of quaternionic Shimura varieties.~I.
\bibitem[JL]{JL} H.~Jacquet, R.~P.~Langlands, \emph{Automorphic forms on $\GL(2)$}, LNM 114, Springer 1970.
\bibitem[KO]{KO} C.-H.~Kim, K.~Ota, On the quantitative variation of congruence ideals and integral periods of modular forms, Res.\ Math.\ Sci.\ 10 (2023); arXiv:1905.02926.
\bibitem[LM]{LM} J.~Ludwig, C.~Merten, Formalising the Bruhat--Tits tree, Ann.\ Formalized Math.\ 2 (2026), 55--81.
\bibitem[LV]{LV} M.~Longo, S.~Vigni, arXiv:0903.2797 ($\Lambda$-adic JL for Hida families on definite quaternion algebras).
\bibitem[MP]{MP} I.~Moerdijk, D.~A.~Pronk, Simplicial cohomology of orbifolds, Indag.\ Math.\ 10 (1999).
\bibitem[MW]{MW} B.~Mazur, A.~Wiles, On $p$-adic analytic families of Galois representations, Compositio Math.\ 59 (1986).
\bibitem[Mz]{Mz} B.~Mazur, Modular curves and the Eisenstein ideal, Publ.\ Math.\ IHÉS 47 (1977).
\bibitem[PW]{PW} R.~Pollack, T.~Weston, On anticyclotomic $\mu$-invariants of modular forms, Compositio Math.\ 147 (2011); arXiv:math/0610694.
\bibitem[Ri84]{Ri84} K.~Ribet, Congruence relations between modular forms, Proc.\ ICM Warszawa 1983 (1984), 503--514.
\bibitem[Ri90]{Ri90} K.~Ribet, On modular representations of $\mathrm{Gal}(\bar\Q/\Q)$ arising from modular forms, Invent.\ Math.\ 100 (1990).
\bibitem[RS]{RS} K.~Ribet, W.~Stein, Lectures on Serre's conjectures.
\bibitem[Se]{Se} J.-P.~Serre, \emph{Trees}, Springer.
\bibitem[SGA7]{SGA7} A.~Grothendieck, Modèles de Néron et monodromie, SGA~7~I, exposé IX, LNM 288.
\bibitem[Vo]{Vo} J.~Voight, \emph{Quaternion algebras}, GTM 288, Springer 2021.
\bibitem[Wi95]{Wi95} A.~Wiles, Modular elliptic curves and Fermat's Last Theorem, Ann.\ of Math.\ 141 (1995).
\end{thebibliography}
\end{document}
